\documentclass[11pt,a4paper]{article}
\usepackage[utf8]{inputenc}
\usepackage[T1]{fontenc}
\usepackage[portuguese]{babel}
\usepackage{geometry}
\geometry{margin=2.5cm}
\usepackage{listings}
\usepackage{xcolor}
\usepackage{amsmath}
\usepackage{amssymb}
\usepackage{hyperref}
\usepackage{enumitem}
\usepackage{tcolorbox}
\usepackage{tabularx}
\usepackage{booktabs}

\definecolor{codebg}{rgb}{0.95,0.95,0.95}
\definecolor{codegreen}{rgb}{0.1,0.5,0.1}
\definecolor{codepurple}{rgb}{0.5,0.1,0.5}
\definecolor{codegray}{rgb}{0.5,0.5,0.5}

\lstdefinestyle{python}{
    backgroundcolor=\color{codebg},
    commentstyle=\color{codegray}\itshape,
    keywordstyle=\color{codepurple}\bfseries,
    stringstyle=\color{codegreen},
    basicstyle=\ttfamily\small,
    breaklines=true,
    frame=single,
    language=Python,
    showstringspaces=false,
    tabsize=4,
    literate={á}{{\'a}}1 {ã}{{\~a}}1 {é}{{\'e}}1 {í}{{\'i}}1 {ó}{{\'o}}1 {õ}{{\~o}}1 {ú}{{\'u}}1 {ç}{{\c{c}}}1 {Á}{{\'A}}1 {Ã}{{\~A}}1 {É}{{\'E}}1 {Í}{{\'I}}1 {Ó}{{\'O}}1 {Õ}{{\~O}}1 {Ú}{{\'U}}1 {Ç}{{\c{C}}}1 {â}{{\^a}}1 {ê}{{\^e}}1 {ô}{{\^o}}1 {û}{{\^u}}1 {Â}{{\^A}}1 {Ê}{{\^E}}1 {Ô}{{\^O}}1 {Û}{{\^U}}1 {à}{{\`a}}1 {è}{{\`e}}1 {ì}{{\`i}}1 {ò}{{\`o}}1 {ù}{{\`u}}1 {ä}{{\"a}}1 {ö}{{\"o}}1 {Ä}{{\"A}}1 {Ö}{{\"O}}1
}
\lstset{style=python}

\newtcolorbox{objetivos}[1][]{
  colback=green!5!white,
  colframe=green!40!black,
  title=O que você vai aprender nesta página,
  fonttitle=\bfseries,
  #1
}

\newtcolorbox{exercicio}[1]{
  colback=blue!5!white,
  colframe=blue!40!black,
  title=#1,
  fonttitle=\bfseries
}

\newtcolorbox{solucao}{
  colback=yellow!5!white,
  colframe=orange!60!black,
  title=Solução proposta,
  fonttitle=\bfseries
}

\title{\textbf{Data Analysis with Python 2024--2025}\\Parte 6: Machine Learning Types\\\large Soluções dos Exercícios}
\author{Tradução e Adaptação: University of Helsinki --- MOOC.fi}
\date{}

\begin{document}

\maketitle

\section*{Nota Importante}
Este material é uma tradução e adaptação baseada no curso \textit{Data Analysis with Python 2024-2025}, oferecido pela \textbf{Universidade de Helsinque (MOOC.fi)}. As soluções apresentadas a seguir foram elaboradas para fins didáticos, seguindo os enunciados dos exercícios propostos no curso original.

\tableofcontents
\newpage

\section{Soluções - Capítulo 6 (PCA)}

\subsection*{Exercício 6.8 -- Variância Explicada}

\begin{solucao}
\begin{lstlisting}
import pandas as pd
import numpy as np
import matplotlib.pyplot as plt
from sklearn.decomposition import PCA

def explained_variance():
    # Carrega base bruta de treinamento
    df = pd.read_csv("src/data.tsv", sep='\t')
    
    # 1. Calcula a variancia original das colunas (matematica base)
    var = np.var(df, axis=0).values
    
    # 2. Utiliza PCA como redutor e captura a metrica analitica explicativa
    pca = PCA()
    pca.fit(df)
    exp_var = pca.explained_variance_
    
    return var, exp_var

def main():
    var, exp_var = explained_variance()
    
    # Prepara o formatador string padrao .3f separando por espaco simples via .join e map
    var_str = " ".join(f"{x:.3f}" for x in var)
    exp_var_str = " ".join(f"{x:.3f}" for x in exp_var)
    
    print("The variances are:", var_str)
    print("The explained variances after PCA are:", exp_var_str)
    
    # === PARTE 2: PLOTAGEM ===
    # Cumsum gera a progressao logica. Ex: se variance[1] for 15, e var[2] for 15, ele plotara (15, 30).
    cum_exp_var = np.cumsum(exp_var)
    plt.plot(np.arange(1, len(cum_exp_var) + 1), cum_exp_var)
    # plt.show() # REMOVER DO COMMIT
\end{lstlisting}
\end{solucao}

\end{document}
